\documentclass[11pt]{article}

% Change "review" to "final" to generate the final (sometimes called camera-ready) version.
% Change to "preprint" to generate a non-anonymous version with page numbers.
\usepackage[review]{acl}

% Standard package includes
\usepackage{times}
\usepackage{latexsym}

% For proper rendering and hyphenation of words containing Latin characters (including in bib files)
\usepackage[T1]{fontenc}
% For Vietnamese characters
% \usepackage[T5]{fontenc}
% See https://www.latex-project.org/help/documentation/encguide.pdf for other character sets

% This assumes your files are encoded as UTF8
\usepackage[utf8]{inputenc}

% This is not strictly necessary, and may be commented out,
% but it will improve the layout of the manuscript,
% and will typically save some space.
\usepackage{microtype}

% This is also not strictly necessary, and may be commented out.
% However, it will improve the aesthetics of text in
% the typewriter font.
\usepackage{inconsolata}

%Including images in your LaTeX document requires adding
%additional package(s)
\usepackage{graphicx}
\usepackage{booktabs}
\usepackage{multirow}
\usepackage{amssymb}

% If the title and author information does not fit in the area allocated, uncomment the following
%
%\setlength\titlebox{<dim>}
%
% and set <dim> to something 5cm or larger.
\newcommand{\corpus}{{\tt MemeFal}}
\newcommand{\red}[1]{\textcolor{red}{#1}}
\title{{\corpus}: Detecting Logical Fallacies of Arguments in Memes}

% Author information can be set in various styles:
% For several authors from the same institution:
% \author{Author 1 \and ... \and Author n \\
%         Address line \\ ... \\ Address line}
% if the names do not fit well on one line use
%         Author 1 \\ {\bf Author 2} \\ ... \\ {\bf Author n} \\
% For authors from different institutions:
% \author{Author 1 \\ Address line \\  ... \\ Address line
%         \And  ... \And
%         Author n \\ Address line \\ ... \\ Address line}
% To start a separate ``row'' of authors use \AND, as in
% \author{Author 1 \\ Address line \\  ... \\ Address line
%         \AND
%         Author 2 \\ Address line \\ ... \\ Address line \And
%         Author 3 \\ Address line \\ ... \\ Address line}

\author{First Author \\
  Affiliation / Address line 1 \\
  Affiliation / Address line 2 \\
  Affiliation / Address line 3 \\
  \texttt{email@domain} \\\And
  Second Author \\
  Affiliation / Address line 1 \\
  Affiliation / Address line 2 \\
  Affiliation / Address line 3 \\
  \texttt{email@domain} \\}

%\author{
%  \textbf{First Author\textsuperscript{1}},
%  \textbf{Second Author\textsuperscript{1,2}},
%  \textbf{Third T. Author\textsuperscript{1}},
%  \textbf{Fourth Author\textsuperscript{1}},
%\\
%  \textbf{Fifth Author\textsuperscript{1,2}},
%  \textbf{Sixth Author\textsuperscript{1}},
%  \textbf{Seventh Author\textsuperscript{1}},
%  \textbf{Eighth Author \textsuperscript{1,2,3,4}},
%\\
%  \textbf{Ninth Author\textsuperscript{1}},
%  \textbf{Tenth Author\textsuperscript{1}},
%  \textbf{Eleventh E. Author\textsuperscript{1,2,3,4,5}},
%  \textbf{Twelfth Author\textsuperscript{1}},
%\\
%  \textbf{Thirteenth Author\textsuperscript{3}},
%  \textbf{Fourteenth F. Author\textsuperscript{2,4}},
%  \textbf{Fifteenth Author\textsuperscript{1}},
%  \textbf{Sixteenth Author\textsuperscript{1}},
%\\
%  \textbf{Seventeenth S. Author\textsuperscript{4,5}},
%  \textbf{Eighteenth Author\textsuperscript{3,4}},
%  \textbf{Nineteenth N. Author\textsuperscript{2,5}},
%  \textbf{Twentieth Author\textsuperscript{1}}
%\\
%\\
%  \textsuperscript{1}Affiliation 1,
%  \textsuperscript{2}Affiliation 2,
%  \textsuperscript{3}Affiliation 3,
%  \textsuperscript{4}Affiliation 4,
%  \textsuperscript{5}Affiliation 5
%\\
%  \small{
%    \textbf{Correspondence:} \href{mailto:email@domain}{email@domain}
%  }
%}

\begin{document}
\maketitle
\begin{abstract}
Towards AI systems with advanced reasoning, informal logic processing is important. How to enable them to know if an argument is sound and reasonable? 
% When there are errors in the reasoning, we call it logical fallacies.
This paper proposes the task of detecting logical fallacies in vision-language arguments: Given an argument for a claim using multiple premises, detect its logical fallacies if any.
To enable explainability and approach verifiability of labels, we require systems to fill in a logical form of a fallacy whenever they choose it.
Because there has been no prior work on vision-language logical fallacies, we collected {\corpus}, the \textit{first} dataset for the task, comprising [1000] meme-based arguments annotated with the committed fallacies and their logical forms.
Evaluating 15 open-source language models and 2 close source ones on \corpus\ shows that current systems are surprisingly good at [detecting appeal to emotion] but are pool at logos-based fallacies which involve being sharp about logic. {\corpus}, therefore, will serve as a testbed for informal logical reasoning of future AI systems in multimodal settings.
\end{abstract}

%%%%%%%%%%%%%%%% INTRODUCTION %%%%%%%%%%%%%%%%%
\section{Introduction}

Write last



%%%%%%%%%%%%%%%% RELATED WORK %%%%%%%%%%%%%%%%%

\section{Related Work}

\red{Depends on the main focus on the paper, we will expand either the dataset part or the methods part.}

% Start with textual logical fallacies. Cite those 11 papers and their taxonomies. Also cite their way of doing annotations and quality control.
\paragraph{Fallacy Datasets.}
Table~\ref{tab:datasets} summarizes existing fallacy datasets; we describe the most relevant ones here.
Taxonomies vary widely---from single-type studies \citep{habernal_before_2018} to unified benchmarks spanning 23 types \citep{helwe_mafalda_2024}---reflecting the lack of community consensus on fallacy categorization.
Collection methods range from gamification \citep{habernal_argotario_2017} and expert annotation of political debates \citep{goffredo_argument-based_2023} to LLM-assisted crowdsourcing \citep{yeh_cocolofa_2024} and synthetic generation \citep{li_reason_2024}.
A consistent finding is that inter-annotator agreement is low regardless of annotator expertise: \citet{helwe_mafalda_2024} report human F1$\approx$35\% even among PhD-level NLP researchers.
The only prior multimodal work is SemEval-2021 Task 6 \citep{dimitrov_semeval-2021_2021}, which targets persuasion techniques in memes and news images as a multilabel problem---a related but distinct task that does not require logical form annotation.
\corpus{} is the first dataset to focus specifically on logical fallacies in vision-language arguments.

\begin{table}[t]
\centering
\footnotesize
\setlength{\tabcolsep}{4pt}
\begin{tabular}{lrrll}
\toprule
\textbf{Dataset} & \textbf{Size} & \textbf{Types} & \textbf{Domain} & \textbf{Mod.} \\
\midrule
Argotario \citeyearpar{habernal_argotario_2017}        & 1,338  &  5 & Online debates    & T \\
Ad Hominem \citeyearpar{habernal_before_2018}          & 2,085  &  1 & Reddit CMV        & T \\
Reddit Fall. \citeyearpar{sahai_breaking_2021}         & 3,358  &  8 & Reddit            & T \\
SemEval-21 T6 \citeyearpar{dimitrov_semeval-2021_2021} & 7,000+ & 22 & Memes/news        & T+I \\
LOGIC \citeyearpar{jin_logical_2022}                   & 3,528  & 13 & Web/climate       & T \\
ElecDeb \citeyearpar{goffredo_argument-based_2023}     & 2,000  &  6 & Political debates & T \\
MAFALDA \citeyearpar{helwe_mafalda_2024}               & 9,545  & 23 & Mixed             & T \\
CoCoLoFa \citeyearpar{yeh_cocolofa_2024}               & 7,706  &  8 & News comments     & T \\
LFUD \citeyearpar{li_reason_2024}                      & 4,020  & 12 & Synthetic         & T \\
SmartyPat \citeyearpar{xu_socrates_2026}               &   502  & 14 & Reddit            & T \\
FAINA \citeyearpar{ramponi_fine-grained_2025}          & 11,000+& 20 & Social media (it) & T \\
FALCON \citeyearpar{chaves_falcon_2025}                & —      &  6 & COVID tweets      & T \\
\midrule
\corpus{} (ours)                                       & [1000] & 21 & Memes             & T+I \\
\bottomrule
\end{tabular}
\caption{Fallacy datasets. \textbf{Mod.}: T = text, I = image.}
\label{tab:datasets}
\end{table}

% Then talk about methods to detect fallacies: tranformer baselines, structure-aware prompting, neuro-symbolic, finetune for explanation generation.
\paragraph{Fallacy Detection Methods.}
Early supervised methods fine-tune transformers, with \citet{sourati_logical_2023} achieving the strongest pre-LLM result (82.7\% macro F1 on LOGIC) via a prototype-network and case-based reasoning pipeline; however, cross-dataset transfer remains poor \citep{alhindi_multitask_2022}.
Zero-shot LLM prompting \citep{pan_fallacies_2024} outperforms fine-tuned models out-of-distribution but lags in-distribution; structure-aware prompts that prepend logical trees \citep{lei_logical_2024} or contextual elaborations \citep{jeong_ceg_2025} close much of this gap.
Neuro-symbolic approaches couple neural models with formal verifiers---FOL + SMT solving \citep{lalwani_nl2fol_2025} or Prolog-backed decomposition \citep{wang_aid-lf_2025}---catching fallacies that purely neural models miss.
Finally, fine-tuning on structured fallacy explanations improves both detection and downstream reasoning \citep{li_reason_2024,hasanain_propxplain_2025}.

% Move on to multimodal fallacies datasets. Then talk about whether approaches here differ.



%%%%%%%%%%%%%%%% TASK %%%%%%%%%%%%%%%%%



\section{Task Formulation}

\subsection{Vision-Language Fallacies and Their Logical Forms}
\red{Explain why choosing \citet{helwe_mafalda_2024}. Describe taxonomy change from MAFALDA. Emphasize the annotation of logical form to help with explanability and verifiability.}

\subsection{Vision-Language Fallacy Detection (VLFD)}
% \red{Explain the input at argument level}

Let $W$ be the space of token sequences in the application language (English in this paper), $F$ be the set of 20 possible fallacies described above.
% , and $\kappa: F \times \mathbb{N} \rightarrow \mathbb{N}$ be the function that a fallacy and an logical form index to the logical form's number of masks.

Given an argumentative image $I$, the problem is to generate a list of quadruplets $(c_i, P_i, f_i, E_i)$ where $c_i \in W$ is a claim made by the image, $P_i = \{p_{i,1}, p_{i,2}, \ldots, p_{i,k_i}\}$ is the set of premises used by the image to support the claim $c_i$ ($p_{ij} \in W$), $f_i \in F$, and $E_i = \{e_{i,1}, e_{i,2}, \ldots, e_{i,l_i}$ be the set of explanations for the occurences of $f_i$ in the argument $(c_i, P_i)$ ($e_{i, j} \in W$).
% $M_i = \{m_{i,1}, m_{i,2}, \ldots, m_{i,\kappa(f_i)}\}$ are the masks for the .

Constraints:
\begin{itemize}
  \item {\bf Faithfulness.} All arguments ($c_i, P_i$) must be faithful to $I$. That means $c_i$ must be one of the claims being made in the meme, and $P_i$ must be the set of premises used to support $c_i$. 
  \item {\bf Logical form adherence.} All explanations $e_{i, j}$ must follow the template of at least one logical form of $f_i$.
  \item {\bf Multiple fallacies.} An argument can appear multiple times in the input if and only if commit multiple fallacies.
\end{itemize}


%%%%%%%%%%%%%%%% DATASET %%%%%%%%%%%%%%%%%

\section{Annotations}

\subsection{Annotation Procedure}

Fallacy annotation is a cognitively demanding task, where current advanced AI models can significant help by providing a {\it starting point}. For each argument, we used Gemini 3 Flash Preview with temperature=0.2 to pre-generate fallacy annotations. Then, two annotators annotated it across 21 fallacies. For every argument-fallacy instance where disagreement occurs, an expert breaks the tie by providing the third annotation. The gold label will then be the majority label.

\red{Insert prompt for model, guidelines and interface for annotators, guidelines and interface for expert reviewer.}

\subsection{Dataset Statistics and Analysis}

%%%%%%%%%%%%%%%% EVALUATION %%%%%%%%%%%%%%%%%

\section{Evaluation}

\subsection{Metrics}

\subsection{Settings}

\subsection{Results and Analysis}

%%%%%%%%%%%%%%%% CONCLUSION %%%%%%%%%%%%%%%%%
\section{Conclusion}

%%%%%%%%%%%%%%%% LIMITATIONS %%%%%%%%%%%%%%%%%
\section*{Limitations}


% Custom bibliography entries only
\bibliography{custom}

\appendix

\section{Example Appendix}
\label{sec:appendix}

This is an appendix.

\end{document}
